\section{Valeur ajoutée de votre travail}

\begin{guideline}
Cette section décrit ce qui restera après votre départ. Elle doit
permettre à l'évaluateur de mesurer concrètement votre contribution.
Qualifiez et quantifiez autant que possible. Distinguez clairement
votre apport personnel de celui de votre équipe : évitez le « nous »
sans précision.
\end{guideline}


\subsection{Livrables produits}

\begin{guideline}
Listez et décrivez les livrables concrets que vous avez produits :
code, documentation, outils déployés, procédures, rapports, formations
dispensées... Pour chaque livrable, précisez son état à votre arrivée
et son état à votre départ.
\end{guideline}

\begin{center}
\renewcommand{\arraystretch}{1.5}
\begin{tabularx}{\textwidth}{@{}l X X@{}}
\toprule
\textbf{Livrable} & \textbf{État à l'arrivée} & \textbf{État au départ} \\
\midrule
\instruction{Ex. : Rôle Ansible ssh\_hardening} & \instruction{Inexistant} & \instruction{Déployé et validé sur 2 VMs} \\
\instruction{Ajoutez vos livrables...} & \instruction{...} & \instruction{...} \\
\bottomrule
\end{tabularx}
\end{center}


\subsection{Contribution personnelle}

\begin{guideline}
Distinguez précisément ce que vous avez fait vous-même de ce qui a été
fait en équipe. Utilisez le « je » et soyez factuel. Si possible,
quantifiez : temps gagné, taux de couverture, nombre de systèmes
automatisés, réduction d'erreurs...
\end{guideline}

\instruction{[Décrivez votre contribution personnelle de façon précise
et honnête. Ex. : « J'ai conçu et implémenté seul le rôle Ansible de
durcissement SSH, couvrant X paramètres de sécurité et déployé sur
Y machines. » Quantifiez là où c'est possible et référencez les
preuves en annexe.]}


\subsection{Pérennité des réalisations}

\begin{guideline}
Expliquez dans quelle mesure vos réalisations sont réutilisables,
maintenables et documentées. Une contribution qui peut être reprise
et étendue par l'équipe après votre départ a plus de valeur qu'un
livrable one-shot.
\end{guideline}

\instruction{[Expliquez ce qui garantit la pérennité de votre travail :
documentation produite, tests mis en place, modularité du code,
intégration dans les outils existants de l'équipe (Git, AWX, Nexus...),
transmission à un collègue ou successeur...]}